\subsection{出次数2の共有支持の排除}
\begin{theorem}\label{thm:cubic-CY-four-point-degree-two-empty}
四単純基点型の零交換子CY包絡の共有支持に出次数2の頂点があれば、H3は存在しない。
\end{theorem}
\begin{proof}
補題\ref{lem:cubic-CY-four-point-support}により
$x_1=x_4=1,x_3=0$ に規格化する。他の頂点の出次数には条件を課さない。
この規格化を保つ $\operatorname{diag}(1,1,z,1)$ は
$(x_0,x_2,x_5)$ を $(zx_0,zx_2,x_5/z)$ へ移す。

固定出典\cite{AS4CYpinned}の隣接共有係数の五つの結果を、以下の必要条件に限って用いる。
\begin{itemize}
\item $x_0=0,x_2=0$ では、二対角零の角の必要被覆AZ1--AZ3により
$L^\eta_\omega$ 又は $J^\eta_h$ となる。$x_0=0,x_2\ne0$ では
$x_2=1$ に規格化し、第一対角零の必要被覆AF1--AF11により残るのは $J^\eta_h$ である。
\item $x_0\ne0,x_5=0$ では $x_0=1$ とする。第五対角零の必要被覆AV1以下により
残るのは $J^\eta_h$ である。ここでは $x_2$ は自由に保つ。
\item $x_0x_5\ne0$ では $x_0=1$ とし、$p=x_2,f=x_5$ とおく。
$d=(p-1)f-1\ne0$ なら、開部分の必要被覆AO1--AO2の両符号の表示は
三次商次元が全許容値で21以上であるため排除される。
$d=0$ なら $q=p-1\ne0$ として $f=q^{-1}$。
境界の全Schur解の被覆AB1--AB3により、全候補は $T^\eta_\lambda$ と $S_\zeta$ である。
\end{itemize}
$L$ と $J$ は命題\ref{prop:cubic-adjacent-L}と命題\ref{prop:cubic-boundary-J}、
$S$ と $T$ は命題\ref{prop:cubic-adjacent-S}と定理\ref{thm:cubic-adjacent-T-all}で排除される。
従って上の被覆に含まれる全候補が消える。
\end{proof}

共有支持のうち出次数1を含むものと全零支持も既に排除した。
従って零交換子CY領域に残る必要支持は、部分活性14支持・3軌道と全共有1支持・1軌道である。
ここでの軌道は有向共有支持の置換軌道を数えている。

\begin{corollary}\label{cor:cubic-CY-four-point-full-shared-necessary}
零交換子CY包絡を持つH3では、対称射影の基点は四単純で、十二共有係数が全て非零である。
\end{corollary}
\begin{proof}
残る部分活性支持では、出次数は0又は3であり、両方が現れる。
活性頂点を $a$ に移し、独立な対角比で $x_1=x_3=x_4=1$ とする。
固定出典\cite{AS4CYpinned}の部分活性排除（IV1--IV13）は、
この規格化で活性頂点が一つ、二つ、三つの各場合を扱う。
一つの場合は残余式が非零単元を含む。二つの場合は
$r=-x_{10}x_9+x_{13}x_6$ とおくと $x_9r=r(x_{10}r+x_{13}^2)=0$ であり、
$r\ne0$ は $x_{13}\ne0$ と矛盾し、$r=0$ でも残余式に非零単元が現れる。
三つの場合は同出典のIV4--IV13の全局所化分岐を適用する。
この分岐では正則性を仮定せず元Schur解を排除しており、$\lambda\ne0,-1$ の全値が含まれる。
従って部分活性支持は残らない。
\end{proof}

\begin{theorem}\label{thm:cubic-CY-zero-commutator-empty}
H3のCY包絡 $B$ の二次関係空間を $R_B$ とすると、$\dim(R_B\cap\bigwedge^2 B_1)\ge1$ である。
\end{theorem}
\begin{proof}
零交換子と仮定する。系\ref{cor:cubic-CY-four-point-full-shared-necessary}により、
四単純基点で十二共有係数が全て非零である。
各頂点に接する三つの生成元対から得る $B$ の主ブロックを考える。

ある主ブロックが可逆なら、その頂点を $a$ に移し、
$x_1=x_3=x_4=1$ と規格化する。固定出典\cite{AS4CYpinned}の
可逆主ブロックの必要帰着FT1--FT12は、元Schur解をFT2--FT3の
$\mathcal P^\eta_{\lambda,r}$ に被覆する。
この表示の全許容値は定理\ref{thm:cubic-full-P-excluded}で排除される。

全主ブロックが特異であれば、非零の共有係数により階数は1又は2である。
固定出典の特異階数1排除では、階数1の頂点が二つある枝と、
一つの頂点から始める残余分岐を扱う。途中で現れる境界表示
$\mathcal B^\eta_\lambda$ は五次商次元が少なくとも58であり、
必要値56に反する。従って階数1の頂点はない。
全て階数2の枝は同出典の特異階数2排除により、元Schur解が存在しない。
これで矛盾を得る。
\end{proof}
